\section{Realised rebalancing and chronological outcomes}
The aggregate alignment measures do not describe the size of the trades made along each strategy's own path. Table~\ref{tab:trades} reports average monthly L1 turnover together with annualised arithmetic net return and maximum drawdown. Turnover is expressed as a percentage of portfolio weight, not as the fraction of one-way purchases. The return annualisation multiplies the monthly arithmetic mean by twelve; it is not a compound annual growth rate.
\begin{table}[ht]
\centering\small
\caption{Own-path trading and return summaries, 1990--2025. Turnover is mean monthly L1 turnover (\%); return is annualised arithmetic net return (\%); drawdown is the largest wealth decline (\%).}
\label{tab:trades}
\begin{tabular}{lrrr}\toprule
Method & Turnover & Return & Drawdown\\\midrule
Target / joint & 2.29 & 10.07 & -41.89 \\
Raw & 2.04 & 10.01 & -41.02 \\
Scalar & 2.15 & 10.07 & -42.75 \\
Risk only & 2.11 & 10.08 & -42.44 \\
Equal weight & 2.54 & 11.64 & -48.09 \\
Minimum variance & 0.86 & 10.32 & -35.42 \\
Log + variance & 2.18 & 9.99 & -40.72 \\
\bottomrule\end{tabular}
\end{table}
A smaller fitted reward gap does not order these path statistics. The target and joint rewards implement the same fitted preference, but comparisons with raw or scalar rewards also reflect different holdings carried from previous months. The first-order cost convention charges the same coefficient to all realised paths, including the path generated by a reward whose \emph{effective} cost preference differs. This keeps realised accounting distinct from the incentive encoded in the reward.

Table~\ref{tab:periods} breaks the existing paths into three chronological intervals. No strategy is restarted at a boundary, and the estimation window continues to roll through prior observations. The rows describe sample variation within the same experiment; they are not independent training and test replications. The paired primary uncertainty calculation uses the full period, so no separate significance claim is attached to these shorter summaries.
\begin{table}[ht]
\centering\small
\caption{Chronological own-path outcomes. MV is monthly basis points; Sharpe uses monthly net excess returns and $\sqrt{12}$ annualisation.}
\label{tab:periods}
\begin{tabular}{llrr}\toprule
Period & Method & MV & Sharpe\\\midrule
1990--2007 & Target / joint & -14.82 & 0.455 \\
1990--2007 & Raw & -12.98 & 0.460 \\
1990--2007 & Minimum variance & 0.84 & 0.570 \\
2008--2015 & Target / joint & -11.45 & 0.590 \\
2008--2015 & Raw & -9.95 & 0.598 \\
2008--2015 & Minimum variance & 15.07 & 0.778 \\
2016--2025 & Target / joint & 4.53 & 0.735 \\
2016--2025 & Raw & 4.75 & 0.730 \\
2016--2025 & Minimum variance & 4.63 & 0.659 \\
\bottomrule\end{tabular}
\end{table}
These records allow a reader to separate three mechanisms: a reward can alter the preferred portfolio on the current state, the altered portfolio can change future turnover needs, and the market can reward or penalise the resulting exposure. The calibration theorem resolves the first mechanism under its fitted assumptions. The latter two require observed path evidence and a model of estimation and trading conditions. Keeping these mechanisms separate prevents a favourable full-sample statistic from being used as proof of correct reward semantics.
